\chapter{Parsing and Building Messages}

% ============================================================================
\section{Parsing messages}

\subsection{Introduction} % --------------------------------------------------

Parsing is the process of creating a structured representation (for example,
a hierarchy of C++ objects) of a message from its ``textual'' representation
(the raw data that is actually sent on the Internet).

For example, say you have the following email in a file called "hello.eml":

\begin{verbatim}
Date: Thu, Oct 13 2005 15:22:46 +0200
From: Vincent <vincent@vmime.org>
To: you@vmime.org
Subject: Hello from VMime!

A simple message to test VMime
\end{verbatim}

The following code snippet shows how you can easily obtain a
{\vcode vmime::message} object from data in this file:

\begin{lstlisting}[caption={Parsing a message from a file}]
// Read data from file
std::ifstream file;
file.open("hello.eml", std::ios::in | std::ios::binary);

vmime::utility::inputStreamAdapter is(file);

vmime::string data;
vmime::utility::outputStreamStringAdapter os(data);

vmime::utility::bufferedStreamCopy(is, os);

// Actually parse the message
vmime::shared_ptr <vmime::message> msg = vmime::make_shared <vmime::message>();
msg->parse(data);

vmime::shared_ptr <vmime::header> hdr = msg->getHeader();
vmime::shared_ptr <vmime::body> bdy = msg->getBody();

// Now, you can extract some of its components
vmime::charset ch(vmime::charsets::UTF_8);

std::cout
   << "The subject of the message is: "
   << hdr->Subject()->getValue <vmime::text>()->getConvertedText(ch)
   << std::endl
   << "It was sent by: "
   << hdr->From()->getValue <vmime::mailbox>()->getName().getConvertedText(ch)
   << " (email: " << hdr->From()->getValue <vmime::mailbox>()->getEmail() << ")"
   << std::endl;
\end{lstlisting}

The output of this program is:

\begin{verbatim}
The subject of the message is: Hello from VMime!
It was sent by: Vincent (email: vincent@vmime.org)
\end{verbatim}


\subsection{Using the {\vcode vmime::messageParser} object} % ----------------

The {\vcode vmime::messageParser} object allows to parse messages in a more
simple manner. You can obtain all the text parts and attachments as well as
basic fields (expeditor, recipients, subject...), without dealing with
MIME message structure.

\begin{lstlisting}[caption={Using {\vcode vmime::messageParser} to parse
more complex messages}]
// Read data from file
std::ifstream file;
file.open("hello.eml", std::ios::in | std::ios::binary);

vmime::utility::inputStreamAdapter is(file);

vmime::string data;
vmime::utility::outputStreamStringAdapter os(data);

vmime::utility::bufferedStreamCopy(is, os);

// Actually parse the message
vmime::shared_ptr <vmime::message> msg = vmime::make_shared <vmime::message>();
msg->parse(data);

// Here start the differences with the previous example
vmime::messageParser mp(msg);

// Output information about attachments
std::cout << "Message has " << mp.getAttachmentCount()
   << " attachment(s)" << std::endl;

for (int i = 0 ; i < mp.getAttachmentCount() ; ++i) {

   vmime::shared_ptr <const vmime::attachment> att = mp.getAttachmentAt(i);
   std::cout << "   - " << att->getType().generate() << std::endl;
}

// Output information about text parts
std::cout << "Message has " << mp.getTextPartCount()
   << " text part(s)" << std::endl;

for (int i = 0 ; i < mp.getTextPartCount() ; ++i) {

   vmime::shared_ptr <const vmime::textPart> tp = mp.getTextPartAt(i);

   // text/html
   if (tp->getType().getSubType() == vmime::mediaTypes::TEXT_HTML) {

      vmime::shared_ptr <const vmime::htmlTextPart> htp =
         vmime::dynamicCast <const vmime::htmlTextPart>(tp);

      // HTML text is in tp->getText()
      // Plain text is in tp->getPlainText()

      // Enumerate embedded objects
      for (int j = 0 ; j < htp->getObjectCount() ; ++j) {

         vmime::shared_ptr <const vmime::htmlTextPart::embeddedObject> obj =
            htp->getObjectAt(j);

         // Identifier (Content-Id or Content-Location) is obj->getId()
         // Object data is in obj->getData()
      }

   // text/plain or anything else
   } else {

      // Text is in tp->getText()
   }
}
\end{lstlisting}

To get the embedded object referenced by a URL of the HTML text (in the
{\vcode src} attribute of an {\vcode img} element, for example), pass this URL
to {\vcode findObject()}, as it is written in the HTML text, but with its
character references decoded ({\vcode \&amp;} is {\vcode \&}). A
{\vcode cid:} URL is compared with the {\vcode Content-Id} field of objects.
Another URL is compared with their {\vcode Content-Location} field, once
relative URLs are resolved as specified in RFC-2557: against the
{\vcode base} element of the HTML text, else the {\vcode Content-Location}
field of the HTML part or of a part which contains it.

\begin{lstlisting}
// With "Content-Location: http://www.vmime.org/" in the header of the
// message, this finds the object with the following field:
//    Content-Location: http://www.vmime.org/images/logo.png
vmime::shared_ptr <const vmime::htmlTextPart::embeddedObject> obj =
   htp->findObject("images/logo.png");
\end{lstlisting}

Following RFC-2557, objects referenced by a URL are in the
{\vcode multipart/related} part of the HTML text, or in a
{\vcode multipart/related} part which contains it. All of them are listed,
whether or not the HTML text uses them.


\subsection{Non-ASCII text in header fields} % --------------------------------

Non-ASCII text in header fields is usually encoded (RFC 2047 encoded-words,
RFC 2231 parameter values), and the encoding tells its charset. Some messages
contain unencoded non-ASCII text, though: internationalized messages, in which
it is UTF-8\footnote{See RFC 6532: \url{https://www.rfc-editor.org/rfc/rfc6532}},
but also messages from mailers which do not follow the standards, and write it
in another charset.

VMime interprets such text as UTF-8. If you have to parse messages in which it
is not valid UTF-8, you can set the charset to use instead in the parsing
context:

\begin{lstlisting}
vmime::parsingContext ctx;
ctx.setHeaderFallbackCharset(vmime::charset("iso-8859-1"));

vmime::shared_ptr <vmime::message> msg = vmime::make_shared <vmime::message>();
msg->parse(ctx, data);
\end{lstlisting}

Valid UTF-8 text is still interpreted as UTF-8, so internationalized messages
are parsed correctly whatever the fallback charset. This option only applies
to header fields: the charset of a text body is given by its
{\vcode Content-Type} field. You can also set it in the default parsing context
({\vcode vmime::parsingContext::getDefaultContext()}), which is used when no
context is given.


% ============================================================================
\section{Building messages}

\subsection{A simple message\label{msg-building-simple-message}} % -----------

Of course, you can build a MIME message from scratch by creating the various
objects that compose it (parts, fields, etc.). The following is an example of
how to achieve it:

\begin{lstlisting}[caption={Building a simple message from scratch}]
vmime::shared_ptr <vmime::message> msg = vmime::make_shared <vmime::message>();

vmime::shared_ptr <vmime::header> hdr = msg->getHeader();
vmime::shared_ptr <vmime::body> bdy = msg->getBody();

vmime::shared_ptr <vmime::headerFieldFactory> hfFactory =
   vmime::headerFieldFactory::getInstance();

// Append a 'Date:' field
vmime::shared_ptr <vmime::headerField> dateField =
   hfFactory->create(vmime::fields::DATE);

dateField->setValue(vmime::datetime::now());
hdr->appendField(dateField);

// Append a 'Subject:' field
vmime::shared_ptr <vmime::headerField> subjectField =
   hfFactory->create(vmime::fields::SUBJECT);

subjectField->setValue(vmime::text("Message subject"));
hdr->appendField(subjectField);

// Append a 'From:' field
vmime::shared_ptr <vmime::headerField> fromField =
   hfFactory->create(vmime::fields::FROM);

fromField->setValue(vmime::make_shared <vmime::mailbox>("me@vmime.org"));
hdr->appendField(fromField);

// Append a 'To:' field
vmime::shared_ptr <vmime::headerField> toField =
   hfFactory->create(vmime::fields::TO);

vmime::shared_ptr <vmime::mailboxList> recipients =
   vmime::make_shared <vmime::mailboxList>();

recipients->appendMailbox
   (vmime::make_shared <vmime::mailbox>("you@vmime.org"));

toField->setValue(recipients);
hdr->appendField(toField);

// Set the body contents
bdy->setContents(
   vmime::make_shared <vmime::stringContentHandler>(
      "This is the text of your message..."
   )
);

// Output raw message data to standard output
vmime::utility::outputStreamAdapter out(std::cout);
msg->generate(out);
\end{lstlisting}

As you can see, this is a little fastidious. Thankfully, VMime also offers a
more simple way for creating messages. The {\vcode vmime::messageBuilder}
object can create basic messages that you can then customize.

The following code can be used to build exactly the same message as in the
previous example, using the {\vcode vmime::messageBuilder} object:

\begin{lstlisting}[caption={Building a simple message
using {\vcode vmime::messageBuilder}}]
try {

   vmime::messageBuilder mb;

   // Fill in some header fields and message body
   mb.setSubject(vmime::text("Message subject"));
   mb.setExpeditor(vmime::mailbox("me@vmime.org"));
   mb.getRecipients().appendAddress(
      vmime::make_shared <vmime::mailbox>("you@vmime.org")
   );

   mb.getTextPart()->setCharset(vmime::charsets::ISO8859_15);
   mb.getTextPart()->setText(
      vmime::make_shared <vmime::stringContentHandler>(
         "This is the text of your message..."
      )
   );

   // Message construction
   vmime::shared_ptr <vmime::message> msg = mb.construct();

   // Output raw message data to standard output
   vmime::utility::outputStreamAdapter out(std::cout);
   msg->generate(out);

// VMime exception
} catch (vmime::exception& e) {

   std::cerr << "vmime::exception: " << e.what() << std::endl;

// Standard exception
} catch (std::exception& e) {

   std::cerr << "std::exception: " << e.what() << std::endl;
}
\end{lstlisting}

Values which are written as is in header fields, like domain names, media
types or message identifiers, cannot contain line breaks, which would end the
field and start another one, or the body: when generating a header field,
{\vcode vmime::exceptions::invalid\_argument} is thrown if it contains a CR or
LF which is not part of a fold (RFC 5322, section 2.2.3). Text, like the
subject or display names, is encoded instead.


\subsection{Internationalized domain names} % --------------------------------

Domain names can contain non-ASCII characters, as specified by
IDNA2008\footnote{See RFC 5890: \url{https://www.rfc-editor.org/rfc/rfc5890}}.
In header fields, they are usually written in ASCII, as A-labels
({\vcode xn-{}-...}). VMime gives the domain names of email addresses in
UTF-8: when parsing, {\vcode xn-{}-bcher-kva.example} gives
{\vcode bücher.example}. Only valid A-labels are decoded: a label which only
looks like one is kept as is, so that it cannot be displayed as another
domain name.

When generating, non-ASCII domain names are written as A-labels, unless
internationalized email is enabled in the generation context. They are first
mapped as specified by UTS \#46\footnote{See \url{https://www.unicode.org/reports/tr46/}}
(non-transitional processing), so that {\vcode BÜCHER.example} also gives
{\vcode xn-{}-bcher-kva.example}. When internationalized email is enabled,
they are written in UTF-8, as U-labels, with the same mapping:
{\vcode BÜCHER.example} gives {\vcode bücher.example}. A domain name which is
not valid cannot be converted: it is written as is, in UTF-8.

This needs the GNU Libidn2 library or ICU (see the {\vcode VMIME\_IDNA\_LIB}
build option): without them, domain names are not converted. You can check it
with {\vcode vmime::utility::idna::isSupported()}, and convert domain names
yourself:

\begin{lstlisting}
vmime::string ascii;

if (vmime::utility::idna::toASCII("bücher.example", ascii)) {
   // ascii is "xn--bcher-kva.example"
} else {
   // Not a valid domain name, or no IDNA support
}

// "bücher.example"
vmime::string domain =
   vmime::utility::idna::toUnicode("xn--bcher-kva.example");
\end{lstlisting}


\subsection{Adding an attachment} % ------------------------------------------

Dealing with attachments is quite simple. Add the following code to the
previous example to attach a file to the message:

\begin{lstlisting}[caption={Building a message with an attachment using
{\vcode vmime::messageBuilder}}]
// Create an attachment
vmime::shared_ptr <vmime::fileAttachment> att =
   vmime::make_shared <vmime::fileAttachment>(
      /* full path to file */  "/home/vincent/paris.jpg",
      /* content type */       vmime::mediaType("image/jpeg"),
      /* description */        vmime::text("My holidays in Paris")
   );

// You can also set some infos about the file
att->getFileInfo().setFilename("paris.jpg");
att->getFileInfo().setCreationDate(
   vmime::datetime("30 Apr 2003 14:30:00 +0200")
);

// Add this attachment to the message
mb.appendAttachment(att);
\end{lstlisting}


\subsection{HTML messages and embedded objects} % ----------------------------

VMime also supports aggregate messages, which permits to build MIME messages
containing HTML text and embedded objects (such as images). For more information
about aggregate messages, please read RFC-2557 (\emph{MIME Encapsulation of
Aggregate Documents, such as HTML}).

Creating such messages is quite easy, using the {\vcode vmime::messageBuilder}
object. The following code constructs a message containing text in both plain
and HTML format, and a JPEG image:

\begin{lstlisting}[caption={Building an HTML message with an embedded image
using the {\vcode vmime::messageBuilder}}]
// Fill in some header fields
mb.setSubject(vmime::text("An HTML message"));
mb.setExpeditor(vmime::mailbox("me@vmime.org"));
mb.getRecipients().appendAddress(
   vmime::make_shared <vmime::mailbox>("you@vmime.org")
);

// Set the content-type to "text/html": a text part factory must be
// available for the type you are using. The following code will make
// the message builder construct the two text parts.
mb.constructTextPart(
   vmime::mediaType(
      vmime::mediaTypes::TEXT,
      vmime::mediaTypes::TEXT_HTML
   )
);

// Set contents of the text parts; the message is available in two formats:
// HTML and plain text. The HTML format also includes an embedded image.
vmime::shared_ptr <vmime::htmlTextPart> textPart =
   vmime::dynamicCast <vmime::htmlTextPart>(mb.getTextPart());

// -- Add the JPEG image (the reference identifier of the returned object is
// -- used to reference it in the HTML text: "cid:" followed by the famous
// -- "CID", or "Content-Id").
// -- Note: you can also read data from a file; see the next example.
vmime::shared_ptr <const vmime::htmlTextPart::embeddedObject> obj =
   textPart->addObject("<...image data...>",
      vmime::mediaType(vmime::mediaTypes::IMAGE, vmime::mediaTypes::IMAGE_JPEG));

// -- Set the text
textPart->setCharset(vmime::charsets::ISO8859_15);

textPart->setText(
   vmime::make_shared <vmime::stringContentHandler>(
      "This is the <b>HTML text</b>, and the image:<br/>"
      "<img src=\"" + obj->getReferenceId() + "\"/>"
   )
);

textPart->setPlainText(
   vmime::make_shared <vmime::stringContentHandler>(
      "This is the plain text."
   )
);
\end{lstlisting}

This will create a message having the following structure:

\begin{verbatim}
multipart/alternative
   text/plain
   multipart/related
      text/html
      image/jpeg
\end{verbatim}

You can easily tell VMime to read the embedded object data from a file. The
following code opens the file \emph{/path/to/image.jpg}, connects it to an
input stream, then add an embedded object:

\begin{lstlisting}
vmime::shared_ptr <vmime::utility::fileSystemFactory> fs =
   vmime::platform::getHandler()->getFileSystemFactory();

vmime::shared_ptr <vmime::utility::file> imageFile =
   fs->create(fs->stringToPath("/path/to/image.jpg"));

vmime::shared_ptr <vmime::contentHandler> imageCts =
   vmime::make_shared <vmime::streamContentHandler>(
      imageFile->getFileReader()->getInputStream(),
      imageFile->getLength()
   );

vmime::shared_ptr <const vmime::htmlTextPart::embeddedObject> obj =
   textPart->addObject(
      imageCts,
      vmime::mediaType(
         vmime::mediaTypes::IMAGE,
         vmime::mediaTypes::IMAGE_JPEG
      )
   );
\end{lstlisting}

An embedded object can also be referenced by a URL rather than by a
Content-Id, for example to send a web page with its images. The URL is
written in the {\vcode Content-Location} field of the object, and you use the
same URL in the HTML text. It can be absolute or relative, but it cannot
contain white-spaces: encode them as {\vcode \%20}, both in the URL and in the
HTML text. A long URL is folded as specified in RFC-2557.

\begin{lstlisting}
textPart->addObject(
   imageCts,
   vmime::mediaType(
      vmime::mediaTypes::IMAGE,
      vmime::mediaTypes::IMAGE_JPEG
   ),
   "images/paris.jpg"
);

textPart->setText(
   vmime::make_shared <vmime::stringContentHandler>(
      "<img src=\"images/paris.jpg\"/>"
   )
);
\end{lstlisting}


% ============================================================================
\section{Working with attachments: the attachment helper\label{msg-attachment-helper}}

The {\vcode attachmentHelper} object allows listing all attachments in a
message, as well as adding new attachments, without using the
{\vcode messageParser} and {\vcode messageBuilders} objects. It can work
directly on messages and body parts.

To use it, you do not need any knowledge about how attachment parts should
be organized in a MIME message.

The following code snippet tests if a body part is an attachment, and if so,
extract its contents to the standard output:

\begin{lstlisting}[caption={Testing if a body part is an attachment}]
vmime::shared_ptr <vmime::bodyPart> part;  // suppose we have a body part

if (vmime::attachmentHelper::isBodyPartAnAttachment(part)) {

   // The body part contains an attachment, get it
   vmime::shared_ptr <const vmime::attachment> attach =
      vmime::attachmentHelper::getBodyPartAttachment(part);

   // Extract attachment data to standard output
   vmime::utility::outputStreamAdapter out(std::cout);
   attach->getData()->extract(out);
}
\end{lstlisting}

You can also easily extract all attachments from a message:

\begin{lstlisting}[caption={Extracting all attachments from a message}]
vmime::shared_ptr <vmime::message> msg;  // suppose we have a message

const std::vector <vmime::shared_ptr <const vmime::attachment> > atts =
   vmime::attachmentHelper::findAttachmentsInMessage(msg);
\end{lstlisting}

Finally, the {\vcode attachmentHelper} object can be used to add an
attachment to an existing message, whatever it contains (text parts,
attachments, ...). The algorithm can modify the structure of the
message if needed (eg. add a \emph{multipart/mixed} part if no one
exists in the message). Simply call the {\vcode addAttachment}
function:

\begin{lstlisting}[caption={Adding an attachment to an existing message}]
vmime::shared_ptr <vmime::message> msg;  // suppose we have a message

// Create an attachment
vmime::shared_ptr <vmime::fileAttachment> att =
   vmime::make_shared <vmime::fileAttachment>(
      /* full path to file */  "/home/vincent/paris.jpg",
      /* content type */       vmime::mediaType("image/jpeg"),
      /* description */        vmime::text("My holidays in Paris")
   );

// Attach it to the message
vmime::attachmentHelper::addAttachment(msg, att);
\end{lstlisting}

